\documentclass{tum_report}
\usepackage[T1]{fontenc}
\usepackage{combelow}
\usepackage{booktabs}
\usepackage{microtype}

\counterwithout{table}{section}
\counterwithin{table}{chapter}
\captionsetup[table]{justification=centering}
\hypersetup{pdftitle={Laboratory Work No. 4 - Block Ciphers and DES},
    pdfauthor={Mihail Pancenco},pdfsubject={Cryptography, Task 2.8, group list number 19}}
\lstdefinestyle{reportcode}{
    language=[Sharp]C,basicstyle=\ttfamily\fontsize{9}{11}\selectfont,
    keywordstyle=\bfseries\color{blue!45!black},numbers=none,
    columns=fullflexible,keepspaces=true,showstringspaces=false,breaklines=true,
    frame=single,rulecolor=\color{black!25},backgroundcolor=\color{black!2},
    aboveskip=8pt,belowskip=8pt,xleftmargin=4pt,xrightmargin=4pt
}
\lstdefinestyle{console}{style=reportcode,language={},breakatwhitespace=true}
\newcommand{\hex}[1]{\texttt{0x#1}}
\clubpenalty=10000
\widowpenalty=10000

\begin{document}
\ministry{MINISTRY OF EDUCATION AND RESEARCH\\OF THE REPUBLIC OF MOLDOVA}
\university{\textbf{Technical University of Moldova}}
\faculty{\textbf{Faculty of Computers, Informatics and Microelectronics}}
\department{\textbf{Department of Software Engineering and Automatics}}
\title{Cryptography}{}
\reportyear{2026}

\begin{titlepage}
\centering
\includegraphics[width=0.20\textwidth]{UTM_Logo.png}\par\vspace{0.8cm}
{\bfseries
\textsc{\ministryname}\\[0.15cm]
\universityname\\
\facultyname\\
\departmentname\par}

\vspace{1cm}
\textbf{Mihail Pancenco}\quad\textbf{FAF-243}\par
\vspace{0.9cm}
{\Large\textbf{Report}\par}
\vspace{0.5cm}
{\Large\normalfont\textit{Laboratory Work No. 4}\par}
\vspace{0.35cm}
{\Large\normalfont\textit{Cryptography}\par}
\vspace{0.7cm}
{\huge\normalfont\textit{Block Ciphers and DES}\par}
\vspace{0.4cm}
{\large\normalfont Computing the right half of a DES round\\Task 2.8 -- Group list number 19\par}

\vspace{1cm}
\begin{flushright}
\normalfont
Supervisors:\par
asist. univ. \textbf{Zaica Maia}\par
conf. univ., dr. \textbf{Zgureanu Aureliu}\par
DISA, FCIM, UTM
\end{flushright}
\vfill
\textbf{Chișinău -- \degreeyear}\par
\end{titlepage}

\clearpage
\addtocontents{toc}{\protect\thispagestyle{empty}}
\tableofcontents
\thispagestyle{empty}
\clearpage

\chapter{Task and short theory}
\section{Assigned task}
My number in the group list is $n=19$. The assignment uses $n\bmod11$ to
select the exercise, so $19\bmod11=8$ gives Task~2.8. The task is to compute
the 32-bit right half $R_k$ of a DES round from the previous left half
$L_{k-1}$ and the combined output of the eight S-boxes. The application must
accept manual or randomly generated inputs and display the table and
intermediate results used in the calculation~\cite{ref:assignment}.

\section{Where the calculation fits in DES}
DES processes 64-bit blocks in 16 Feistel rounds. Its 64-bit key representation
includes eight parity bits, leaving 56 bits for the key schedule. An initial
permutation produces the two 32-bit halves $L_0$ and $R_0$. Each round uses a
48-bit subkey $K_k$ and computes
\[
L_k=R_{k-1}, \qquad R_k=L_{k-1}\oplus f(R_{k-1},K_k).
\]
After the last round, the inverse initial permutation is applied to
$R_{16}\|L_{16}$ to obtain the ciphertext block~\cite{ref:fips}.

Inside $f$, expansion $E$ maps $R_{k-1}$ to 48 bits, which are XORed with
$K_k$. The result is split into eight 6-bit blocks $B_1,\ldots,B_8$.
Each S-box uses the outer two bits of its input to select a row and
the inner four bits to select a column, producing four output bits.
Concatenating the eight outputs yields the 32-bit word
$S=S_1(B_1)\|\cdots\|S_8(B_8)$. Task~2.8 starts with this word already supplied,
so the program computes $f=P(S)$ and then $R_k=L_{k-1}\oplus P(S)$.

\section{The DES permutation \texorpdfstring{$P$}{P}}
Read Table~\ref{tab:permutation} left to right, row by row. Its $j$-th entry
is the input position copied to output position $j$. Positions are counted
from the most significant bit, starting at one. Thus the first output
bit is input bit 16, and the last is input bit 25~\cite{ref:fips}.

\begin{table}[H]
\caption{Input bit positions in output order}
\label{tab:permutation}
\centering
\setstretch{1}
\begin{tabular}{rrrr}
\toprule
16 & 7  & 20 & 21 \\
29 & 12 & 28 & 17 \\
1  & 15 & 23 & 26 \\
5  & 18 & 31 & 10 \\
2  & 8  & 24 & 14 \\
32 & 27 & 3  & 9  \\
19 & 13 & 30 & 6  \\
22 & 11 & 4  & 25 \\
\bottomrule
\end{tabular}
\end{table}

\clearpage
\chapter{Implementation}
\section{Console input}
The application is written in C\# and targets .NET 10.
Option \texttt{1} prompts for $L_{k-1}$ and $S$ in hexadecimal.
\texttt{Convert.ToUInt32(input, 16)} accepts both \hex{12345678} and
\texttt{12345678}. An invalid menu choice repeats the menu prompt.
A format error or a value larger than \hex{FFFFFFFF} repeats the
corresponding hexadecimal prompt.

Option \texttt{2} uses \texttt{Random.Next(0, int.MaxValue)} to generate
each input. This limits generated words to \hex{00000000} through
\hex{7FFFFFFE}, so their most significant bit is always zero. Manual input
accepts the full unsigned 32-bit range.

\section{Bit selection and XOR}
\texttt{Program.cs} passes $S$ and $L_{k-1}$ to \texttt{Algorithm}.
The method \texttt{P} stores the permutation as a 32-entry byte array and
prints it before selecting bits. To read input position $p$, the code
shifts the 32-bit unsigned word right by $32-p$ and masks it with
\texttt{1}. It shifts the partial result left and appends the selected bit
using bitwise OR. After 32 iterations the accumulator contains $P(S)$.

{\setstretch{1}
\lstinputlisting[style=reportcode,linerange={68-70,73-79},gobble=8]{../code/task_08/Algorithm.cs}
}

The constructor computes \texttt{this.Rk = this.Ln\_minus\_one \string^ permutated}.
The C\# \texttt{\string^} operator XORs corresponding bits, yielding one where
they differ. The trace shows the inputs, the $P$ table, the permuted word,
and the XOR operands and result. Binary words are padded to 32 bits and
hexadecimal words to eight digits. The main program prints the value
returned by \texttt{getResult()}.

\section{Project entry point}
Run the following command from the repository root.
{\setstretch{1}
\begin{lstlisting}[style=console]
dotnet run --project lab_04/code/DES_cipher.csproj
\end{lstlisting}
}
The console entry point runs the two steps assigned in Task~2.8.

The Task~2.8 entry point takes the S-box output
directly and uses \path{task_08/Algorithm.cs} for the round calculation.

\clearpage
\chapter{Program output}
This run uses option~\texttt{2} to generate both input words.

{\setstretch{1}
\lstinputlisting[style=console]{evidence/random.txt}
}

\clearpage
\chapter{Conclusions}
In this laboratory work, I implemented the last two steps of the DES
round function and update needed to find the new right half. The
assignment starts with the combined output of the eight S-boxes, so
the same calculation can be used for any round once that word and
the previous left half are available.

The permutation required careful attention to bit numbering. The DES
table counts positions from the most significant bit, while the C\#
code uses right shifts to read a bit. Translating a table position
$p$ into a shift of $32-p$ lets the program build the output word in
the order specified by the table.

The XOR then combines the permuted word with the previous left half.
Its reversibility is part of the Feistel structure. Applying the
same XOR to $R_k$ recovers $L_{k-1}$, and the previous right half is
available as $L_k$. This connects the two implemented steps to the
way DES reverses a round.

The random mode provides inputs for a quick demonstration. To repeat
that calculation, I can enter the two displayed hexadecimal values
in manual mode. Showing the binary and hexadecimal forms together
makes it easier to follow the bit operations and read the final
32-bit result.

\par
\vspace{0.6cm}
\phantomsection
\addcontentsline{toc}{chapter}{Bibliography}
\begin{thebibliography}{9}
\singlespace\normalsize\raggedright
\bibitem{ref:assignment}
Technical University of Moldova.
\textit{Lucrare de laborator nr. 4. Cifruri bloc. Algoritmul DES}.
Course assignment, p.~1, Task~2.8.

\bibitem{ref:fips}
National Institute of Standards and Technology.
\textit{Data Encryption Standard (DES)}. FIPS PUB 46-3,
25 October 1999, pp.~8--14 and Appendix~1.
\url{https://csrc.nist.gov/pubs/fips/46-3/final}.
\end{thebibliography}

\end{document}
